%=====================================================================
\chapter{Probability and Uncertainty}\label{ch:probability}
%=====================================================================
\locator{1}{Part I --- Foundations of Data}

\begin{outcomes}
By the end of this chapter you will be able to:
\begin{tight}
  \item \textbf{Apply} the addition and multiplication rules of probability, and test whether two events are independent.
  \item \textbf{Compute} a conditional probability and explain why $P(A \mid B) \neq P(B \mid A)$.
  \item \textbf{Use} Bayes' theorem to update a belief in the light of evidence.
  \item \textbf{Explain} the base-rate fallacy and demonstrate it numerically for a rare-event detector.
  \item \textbf{State} the Central Limit Theorem and say why it makes inference possible.
\end{tight}
\end{outcomes}

\begin{prereq}
\chref{descriptive} --- distributions, mean and standard deviation. Bayes'
theorem here is the direct foundation of naive Bayes in \chref{classification} and
of the false-alarm arithmetic in \chref{evaluation} and \chref{cybersecurity}.
\end{prereq}

\begin{keyterms}
experiment $\bullet$ outcome $\bullet$ event $\bullet$ sample space $\bullet$
independence $\bullet$ mutually exclusive $\bullet$ conditional probability
$\bullet$ Bayes' theorem $\bullet$ prior, likelihood, posterior $\bullet$ base rate
$\bullet$ random variable $\bullet$ expected value $\bullet$ Central Limit Theorem
\end{keyterms}

\section{Why Probability Is the Language of This Course}

Every prediction you make is uncertain, and pretending otherwise is the fastest
way to lose the trust of the people who use your work. Probability is how
uncertainty is written down, compared, and combined. It is also, quite directly,
what a classifier outputs: ``0.78 probability of failing'' is a probability
statement, and you cannot interpret it without this chapter.

\section{The Rules}

\begin{definitionbox}[Probability]
For an event $A$, $P(A)$ is a number in $[0,1]$ giving the long-run proportion of
occasions on which $A$ happens. $P(A)=0$ means impossible, $P(A)=1$ means certain,
and $P(\text{not }A) = 1 - P(A)$.
\end{definitionbox}

\rowcolors{2}{white}{lightbg}
\begin{longtable}{@{}p{3.3cm}p{5.0cm}p{6.3cm}@{}}
\toprule
\rowcolor{primary!15}
\textbf{Rule} & \textbf{Formula} & \textbf{In words} \\
\midrule
\endhead
Complement & $P(\bar{A}) = 1 - P(A)$ & Either it happens or it does not \\
Addition (general) & $P(A \cup B) = P(A) + P(B) - P(A \cap B)$ & Add, then subtract the double-counted overlap \\
Addition (exclusive) & $P(A \cup B) = P(A) + P(B)$ & Valid \emph{only} when $A$ and $B$ cannot both occur \\
Multiplication (general) & $P(A \cap B) = P(A)\,P(B \mid A)$ & Chance of $A$, then chance of $B$ given $A$ happened \\
Multiplication (independent) & $P(A \cap B) = P(A)\,P(B)$ & Valid \emph{only} when $A$ tells you nothing about $B$ \\
\bottomrule
\end{longtable}

\begin{alertbox}[The Two Words That Cause Most Errors]
\emph{Mutually exclusive} and \emph{independent} are not the same thing and are
almost opposites. Mutually exclusive events \emph{cannot} co-occur, so knowing one
happened tells you the other definitely did not --- which is the strongest possible
\emph{dependence}. Two mutually exclusive events with non-zero probability are
never independent.
\end{alertbox}

\section{Conditional Probability}

\begin{definitionbox}[Conditional Probability]
\[
P(A \mid B) \;=\; \frac{P(A \cap B)}{P(B)}, \qquad P(B) > 0
\]
Read as ``the probability of $A$ \emph{given that} $B$ has occurred''. It is the
probability of $A$ within the reduced world in which $B$ is true.
\end{definitionbox}

\dsfig{%
\begin{tikzpicture}[font=\scriptsize]
% --- unconditional world
\draw[draw=gray!55, fill=gray!6, line width=1pt] (0,0) rectangle (5.0,3.2);
\node[anchor=north west, text=gray!55!black] at (0.08,3.15) {all sessions $S$};
\begin{scope}
  \clip (0,0) rectangle (5.0,3.2);
  \fill[secondary!22] (1.75,1.6) circle (1.3cm);
  \fill[accent!28] (3.15,1.6) circle (1.05cm);
  \begin{scope}
    \clip (1.75,1.6) circle (1.3cm);
    \fill[purplehl!45] (3.15,1.6) circle (1.05cm);
  \end{scope}
\end{scope}
\draw[secondary, line width=1pt] (1.75,1.6) circle (1.3cm);
\draw[accent, line width=1pt] (3.15,1.6) circle (1.05cm);
\node[text=secondary!60!black, font=\scriptsize\bfseries] at (1.15,2.55) {$B$: from};
\node[text=secondary!60!black, font=\scriptsize\bfseries] at (1.15,2.28) {new device};
\node[text=accent!70!black, font=\scriptsize\bfseries] at (4.05,2.5) {$A$: fraud};
\node[text=purplehl!50!black, font=\tiny\bfseries] at (2.62,1.6) {$A\cap B$};

\node[align=center, text=gray!55!black] at (2.5,-0.55)
     {$P(A)=\dfrac{|A|}{|S|}$ --- fraud among \emph{all} sessions};

% --- arrow
\draw[-{Stealth[length=3mm]}, primary!60, line width=2pt] (5.6,1.6) -- (7.1,1.6);
\node[text=primary!70!black, font=\scriptsize, align=center] at (6.35,2.15) {condition\\on $B$};

% --- conditioned world
\draw[secondary, line width=1.4pt, fill=secondary!10] (9.9,1.6) circle (1.55cm);
\begin{scope}
  \clip (9.9,1.6) circle (1.55cm);
  \fill[purplehl!45] (10.75,1.6) circle (1.25cm);
  \draw[accent, line width=1pt] (10.75,1.6) circle (1.25cm);
\end{scope}
\node[text=secondary!60!black, font=\scriptsize\bfseries] at (9.9,3.4) {the world shrinks to $B$};
\node[text=purplehl!50!black, font=\tiny\bfseries] at (10.55,1.6) {$A\cap B$};
\node[align=center, text=gray!55!black] at (9.9,-0.55)
     {$P(A\mid B)=\dfrac{|A\cap B|}{|B|}$ --- fraud among \emph{new-device} sessions};
\end{tikzpicture}}%
{Conditioning shrinks the universe. $P(A)$ divides the overlap by everything; $P(A \mid B)$
divides the same overlap by $B$ alone. Because the denominators differ,
$P(A\mid B)$ and $P(B \mid A)$ are different numbers --- confusing them is the base-rate
fallacy.}{conditional-prob}

\section{Bayes' Theorem}

\begin{definitionbox}[Bayes' Theorem]
\[
\underbrace{P(A \mid B)}_{\text{posterior}}
\;=\;
\frac{\overbrace{P(B \mid A)}^{\text{likelihood}} \; \overbrace{P(A)}^{\text{prior}}}
     {\underbrace{P(B)}_{\text{evidence}}}
\]
It converts $P(B \mid A)$ --- which you can usually measure --- into
$P(A \mid B)$ --- which is usually what you actually want to know.
\end{definitionbox}

\begin{worked}[The Base-Rate Fallacy: An Intrusion Detector]
An intrusion detection system is tested and found to be excellent:
\begin{tight}
  \item It flags 99\% of genuine attacks: $P(\text{alert} \mid \text{attack}) = 0.99$.
  \item It has a 1\% false alarm rate: $P(\text{alert} \mid \text{normal}) = 0.01$.
\end{tight}
Attacks are rare: 1 session in 10{,}000 is an attack, so $P(\text{attack}) = 0.0001$.

\smallskip
\textbf{The alert fires. What is the probability this is a real attack?}

\smallskip
Work with a concrete population of 1{,}000{,}000 sessions.
\begin{tight}
  \item Genuine attacks: $1{,}000{,}000 \times 0.0001 = 100$.
        Of these, $100 \times 0.99 = \mathbf{99}$ are flagged.
  \item Normal sessions: $999{,}900$.
        Of these, $999{,}900 \times 0.01 = \mathbf{9{,}999}$ are flagged.
\end{tight}

\textbf{Total alerts:} $99 + 9{,}999 = 10{,}098$.

\[
P(\text{attack} \mid \text{alert}) \;=\; \frac{99}{10{,}098} \;\approx\; 0.0098
\]

\smallskip
\textbf{Under 1\%.} A detector that catches 99\% of attacks produces alerts that
are wrong 99\% of the time. Nothing is broken --- the arithmetic is simply
dominated by the \term{base rate}. There are so many more normal sessions that even
a 1\% error on them swamps the true positives.

\smallskip
\textbf{Consequences you must remember:}
\begin{tight}
  \item This is why security teams suffer alert fatigue, and why reducing the
        false-positive rate matters far more than raising the detection rate.
  \item To reach $P(\text{attack}\mid\text{alert}) > 0.5$ here, the false alarm
        rate must fall to about $0.0001$ --- a hundredfold improvement.
  \item The same arithmetic governs medical screening, plagiarism detection and
        every rare-event problem you will meet. \chref{evaluation} formalises it
        as \emph{precision}.
\end{tight}
\end{worked}

\dsfig{%
\begin{tikzpicture}
\begin{axis}[
  width=12.0cm, height=5.6cm,
  xmode=log, xmin=0.00002, xmax=0.5,
  ymin=0, ymax=1.02,
  xlabel={\scriptsize base rate --- proportion of sessions that really are attacks (log scale)},
  ylabel={\scriptsize $P(\text{attack}\mid\text{alert})$},
  tick label style={font=\tiny}, label style={font=\scriptsize},
  axis lines=left, grid=major, grid style={gray!18},
  legend style={font=\tiny, draw=gray!40, at={(0.02,0.98)}, anchor=north west},
]
% P(A|alert) = .99b / (.99b + fpr(1-b))
\addplot[accent, line width=1.3pt, domain=0.00002:0.5, samples=200]
  {0.99*x/(0.99*x + 0.01*(1-x))};
\addlegendentry{false alarm rate 1\%}
\addplot[goldhl, line width=1.3pt, domain=0.00002:0.5, samples=200]
  {0.99*x/(0.99*x + 0.001*(1-x))};
\addlegendentry{false alarm rate 0.1\%}
\addplot[greenhl, line width=1.3pt, domain=0.00002:0.5, samples=200]
  {0.99*x/(0.99*x + 0.0001*(1-x))};
\addlegendentry{false alarm rate 0.01\%}

\addplot[primary, dashed, line width=0.8pt, domain=0.00002:0.5] {0.5};
\node[font=\tiny, primary, anchor=south east] at (axis cs:0.45,0.51) {alerts more right than wrong};

\addplot[black, only marks, mark=*, mark size=1.8pt] coordinates {(0.0001,0.0098)};
\draw[gray!60, line width=0.5pt] (axis cs:0.0001,0.0098) -- (axis cs:0.00025,0.30);
\node[font=\tiny, anchor=south west, align=left, fill=white, fill opacity=0.85,
      text opacity=1, inner sep=1.5pt] at (axis cs:0.00023,0.30)
     {the worked example:\\1 in 10{,}000, 1\% FPR\\$\Rightarrow$ 0.98\%};
\end{axis}
\end{tikzpicture}}%
{Why rare-event detection is hard. Detection rate is held at 99\% throughout; only the
false alarm rate changes. At a base rate of 1 in 10{,}000 even a 0.1\% false alarm rate
leaves most alerts wrong. This single figure explains the design of every alerting system
in \chref{cybersecurity}.}{base-rate}

\section{Random Variables and Expected Value}

\begin{definitionbox}[Random Variable and Expected Value]
A \term{random variable} $X$ assigns a number to each outcome. Its \term{expected
value} is the long-run average,
\[
E[X] \;=\; \sum_{k} k \cdot P(X = k)
\]
--- each possible value weighted by how often it occurs.
\end{definitionbox}

\begin{worked}[Is the Alert Worth Investigating?]
Each investigation costs 20 minutes of analyst time. A missed attack costs an
estimated 40 hours of incident response. Using the numbers above,
$P(\text{attack} \mid \text{alert}) = 0.0098$.

\smallskip
\textbf{Expected cost of investigating:} 20 minutes, always.

\textbf{Expected cost of ignoring:}
$0.0098 \times 2400 \text{ min} + 0.9902 \times 0 = 23.5$ minutes.

\smallskip
\textbf{Conclusion:} investigate --- but only just. The margin is 3.5 minutes per
alert, so a modest rise in the false alarm rate would flip the decision and make
the detector actively harmful to deploy. Expected value turns ``is this model good
enough?'' from an opinion into an arithmetic question, which is exactly how it
should be presented to the people funding the work.
\end{worked}

\section{The Central Limit Theorem}

\begin{definitionbox}[Central Limit Theorem (CLT)]
Take samples of size $n$ from \emph{any} distribution with a finite mean $\mu$ and
standard deviation $\sigma$. As $n$ grows, the distribution of the \emph{sample
mean} approaches a normal distribution with mean $\mu$ and standard deviation
$\sigma/\sqrt{n}$ --- regardless of the shape of the original distribution.
\end{definitionbox}

\dsfig{%
\begin{tikzpicture}
\begin{groupplot}[
  group style={group size=4 by 1, horizontal sep=0.6cm},
  width=4.35cm, height=3.4cm,
  axis lines=left, ymin=0, ytick=\empty,
  tick label style={font=\tiny}, title style={font=\tiny\bfseries},
]
\nextgroupplot[title={the population ($n{=}1$)}, xmin=0, xmax=6]
\addplot[draw=accent, fill=accent!18, line width=0.9pt, domain=0.03:6, samples=140]
  {1.9*exp(-(ln(x)-0.05)^2/(2*0.8^2))/(x*0.8*2.5066)};
\node[font=\tiny, accent!80!black] at (axis cs:3.6,0.55) {very skewed};

\nextgroupplot[title={means of $n{=}5$}, xmin=0, xmax=6]
\addplot[draw=goldhl, fill=goldhl!22, line width=0.9pt, domain=0.03:6, samples=140]
  {1.55*exp(-(ln(x)-0.25)^2/(2*0.42^2))/(x*0.42*2.5066)};

\nextgroupplot[title={means of $n{=}20$}, xmin=0, xmax=6]
\addplot[draw=secondary, fill=secondary!22, line width=0.9pt, domain=0.03:6, samples=140]
  {1.35*exp(-(ln(x)-0.32)^2/(2*0.22^2))/(x*0.22*2.5066)};

\nextgroupplot[title={means of $n{=}100$}, xmin=0, xmax=6]
\addplot[draw=greenhl, fill=greenhl!22, line width=0.9pt, domain=0.03:6, samples=140]
  {1.28*exp(-(ln(x)-0.35)^2/(2*0.10^2))/(x*0.10*2.5066)};
\node[font=\tiny, greenhl!55!black, align=center] at (axis cs:3.9,1.6) {normal,\\and narrow};
\end{groupplot}
\end{tikzpicture}}%
{The Central Limit Theorem at work. The population (left) is heavily skewed, yet the
distribution of \emph{sample means} becomes normal and progressively narrower as $n$ grows.
The spread shrinks as $\sigma/\sqrt{n}$ --- quadrupling the sample halves the
uncertainty.}{clt}

\begin{conceptbox}[Why the CLT Is the Most Useful Theorem in the Course]
It means you can put a confidence interval around an average, and run a hypothesis
test, \emph{without knowing the shape of the underlying data}. Response times are
wildly skewed; the average of 200 response times is nevertheless approximately
normal, so the standard tools apply. That is the entire licence on which
\chref{inference} operates.

Note also the $\sqrt{n}$: to halve your uncertainty you need \emph{four times} the
data. This is why ``just collect more data'' has sharply diminishing returns, and
why it is worth knowing before you promise a stakeholder anything.
\end{conceptbox}

\begin{pitfall}
\begin{tight}
  \item \textbf{Confusing $P(A\mid B)$ with $P(B\mid A)$.} ``99\% of attacks
        trigger an alert'' is not ``99\% of alerts are attacks''. See
        \dsref{base-rate}.
  \item \textbf{Multiplying probabilities that are not independent.} Two sensors on
        the same power supply do not fail independently, so $P(\text{both fail})
        \gg P(A)P(B)$.
  \item \textbf{Applying the CLT to a single observation.} The theorem is about the
        distribution of \emph{means}, not of raw values. Individual response times
        stay skewed no matter how much data you collect.
\end{tight}
\end{pitfall}

%---------------------------------------------------------------------
\begin{fourlenses}{Conditional Probability and Base Rates}
\lensLA{Only 8\% of students eventually fail. A model that flags at-risk students
with a 10\% false-positive rate will produce a flagged list that is mostly students
who would have passed anyway. Before any outreach programme is designed, compute
$P(\text{fails} \mid \text{flagged})$ --- because a list that is 70\% wrong will
destroy staff confidence in the system within a term, and may harm the students
wrongly labelled (\chref{ethics}).}

\lensPM{The prior probability that a randomly chosen project overruns is about
50\% in many organisations --- which is a \emph{high} base rate, and therefore the
opposite situation to security. Here even a moderately accurate predictor produces
alerts that are usually right, so the practical difficulty is not false alarms but
persuading people to act on a warning they already half-expected.}

\lensIOT{$P(\text{failure within 30 days}) \approx 0.002$ for a given motor. This
is a rare-event problem, so the base-rate arithmetic of \dsref{base-rate} applies
directly: an 80\%-detection, 2\%-false-alarm model yields
$P(\text{fail} \mid \text{alert}) \approx 7\%$. Whether that is acceptable is an
expected-value question --- compare the cost of 13 unnecessary inspections against
one unplanned line stoppage.}

\lensSEC{The worked example is this lens. Add to it: an adversary who knows your
detector exists will deliberately operate just below its threshold, which changes
$P(\text{alert}\mid\text{attack})$ over time in a way no other domain in this course
experiences. Base-rate arithmetic tells you today's precision; \chref{cybersecurity}
deals with the fact that tomorrow's is being actively degraded.}
\end{fourlenses}

%---------------------------------------------------------------------
\begin{lab}[Simulating the Base-Rate Fallacy]
Do not take the arithmetic on trust --- generate a million sessions and count.

\begin{lstlisting}[language=Python]
import numpy as np
rng = np.random.default_rng(42)

N          = 1_000_000
base_rate  = 0.0001     # 1 in 10,000 sessions is an attack
detection  = 0.99       # P(alert | attack)
false_alarm= 0.01       # P(alert | normal)

is_attack = rng.random(N) < base_rate
p_alert   = np.where(is_attack, detection, false_alarm)
alert     = rng.random(N) < p_alert

print("attacks           :", is_attack.sum())
print("alerts raised     :", alert.sum())
print("alerts that were real:", (alert & is_attack).sum())
print("P(attack | alert) :", (alert & is_attack).sum() / alert.sum())
# ~0.0098 -- matches the hand calculation

# Now set false_alarm = 0.0001 and re-run. Precision jumps to ~50%.
# Changing detection from 0.99 to 0.999 barely moves it at all.
\end{lstlisting}

\textbf{The point of the last two lines:} in rare-event detection, effort spent
reducing false alarms is worth vastly more than effort spent raising the detection
rate. That is a counter-intuitive engineering priority, and it falls straight out
of Bayes' theorem.
\end{lab}

%---------------------------------------------------------------------
\begin{checkpoint}
\begin{tightnum}
  \item A test is 95\% accurate for a condition affecting 1 in 1{,}000 people.
        Roughly what fraction of positive tests are correct --- about 95\%, about
        50\%, or about 2\%? Justify without a calculator.
  \item Events $A$ and $B$ are mutually exclusive and both have probability 0.3.
        Are they independent? Show why.
  \item The CLT says the sample mean becomes normal as $n$ grows. Does it say
        individual observations become normal?
\end{tightnum}
\end{checkpoint}

%---------------------------------------------------------------------
\begin{chaptersummary}
\begin{tight}
  \item Probability is how uncertainty is written down; every classifier output is
        a probability statement.
  \item $P(A \cap B) = P(A)P(B)$ only under independence; the addition rule needs
        the overlap subtracted unless events are mutually exclusive.
  \item Conditioning shrinks the sample space, which is why $P(A\mid B)$ and
        $P(B\mid A)$ differ.
  \item Bayes' theorem converts measurable likelihoods into the posterior you
        actually want. Under a low base rate, an excellent detector still produces
        mostly-wrong alerts.
  \item Expected value turns model quality into a cost comparison a stakeholder can
        act on.
  \item The CLT makes sample means normal regardless of the population's shape, and
        uncertainty falls only as $\sqrt{n}$.
\end{tight}
\end{chaptersummary}

\begin{reviewq}
  \item \textbf{(MCQ)} $P(A \cap B) = P(A) \times P(B)$ holds when $A$ and $B$ are:
        (a)~mutually exclusive (b)~independent (c)~equally likely (d)~always.
  \item \textbf{(MCQ)} As sample size quadruples, the standard deviation of the
        sample mean: (a)~quadruples (b)~doubles (c)~halves (d)~is unchanged.
  \item \textbf{(Short)} State Bayes' theorem and name each of its four parts.
  \item \textbf{(Short)} Explain the base-rate fallacy to a non-technical manager
        in three sentences, without using a formula.
  \item \textbf{(Short)} Why are two sensors sharing a power supply not independent,
        and what does that do to your estimate of $P(\text{both fail})$?
  \item \textbf{(Applied)} A plagiarism detector flags 90\% of copied assignments
        and falsely flags 5\% of honest ones. 2\% of assignments are copied. For
        10{,}000 submissions, build the full table of counts and compute
        $P(\text{copied} \mid \text{flagged})$. Comment on whether this detector
        should be used to open disciplinary proceedings automatically.
  \item \textbf{(Applied)} Using expected value, decide whether to investigate
        alerts from a detector with $P(\text{attack}\mid\text{alert}) = 0.004$,
        where investigation costs 30 minutes and a missed attack costs 60 hours.
        Show the calculation and state the break-even precision.
\end{reviewq}
